\section{Methodology}

\subsection{Files Submitted}
Below are the directory structure of the project.
For simplicity, irrelevant or simple configuration files are omitted.
\begin{verbatim}
CSCI2720-Project/
|-- report/                                // Source code for report
|-- data/
|   `-- preprocess.py                      // Data preprocessing script
|-- frontend/
`-- backend/
\end{verbatim}

\subsubsection{Frontend}
\begin{verbatim}
frontend/
 `- src/
    |-- main.jsx                       // Root of the app
    |-- App.jsx                        // Main App component
    |-- index.css                      // Theme CSS config.
    |-- TopNav.jsx                     // Navigation bar
    |-- Home.jsx                       // Home page
    |-- Auth.jsx                       // Authentication page
    |-- Leaderboard.jsx                // Leaderboard page
    |-- event/                         // Event-related files
    |-- location/                      // Location-related files
    |-- UserManager/                   // User Manager related files
    |-- constants/                     // Reuseable constants
    |-- hooks/                         // Reuseable hooks
    |-- public/                        // Assets folder
    |-- lib/                           // Helper functions
    `-- components/                    // Reuseable components
        |-- common-table-toolbar.jsx   // Reuseable table tool-bar
        |-- toggle-favourite.jsx       // Button to toggle favourite
        |-- toggle-like.jsx            // Button to toggle like
        |-- page-shell.jsx             // A "Shell" for padding and title
        |-- mode-toggle.jsx            // Shadcn dark mode toggle
        |-- theme-provider.jsx         // Shadcn dark mode theme provider
        `-- ui/                        // Shadcn UI components
\end{verbatim}

\subsubsection{Backend}
\begin{verbatim}
backend/
  |-- .env                               // Environment files storing secret keys
  |-- init-db.js                         // Script for initializing database
  `-- src/
      |-- index.js                       // Root of backend
      |-- app.js                         // Main App component
      |-- libs/
      |   |-- connect-db.js              // For database connecting
      |   `-- register-account.js        // For registering user accounts
      |-- middleware/
      |   |-- error.js                   // For error handling
      |   |-- logging.js                 // For logging all backend requests
      |   |-- require-auth.js            // Check user auth token
      |   |-- require-recaptcha.js       // Check ReCAPTCHA tokens
      |   `-- timeout.js                 // To simulate slow internet connection
      |-- models/                        // Schema definitions
      `-- routes/                        // All backend routes
\end{verbatim}

\subsection{Data Preprocessing}
Preprocessing the cultural events and venues dataset are simple.
Here we briefly outline the process.
\begin{enumerate}
    \item Get government dataset \verb|venues.xml| and \verb|events.xml| from \url{https://data.gov.hk/en-data/dataset/hk-lcsd-event-event-cultural}
    \item Parse the data with Python \verb|xmltodict| package.
    \item Remove data that missing required fields, e.g. latitude and longitude.
    \item Pick 10 random venues in accordance with the requirements.
    \item Dump the data to JSON for backend initialization.
\end{enumerate}
For details, see \verb|data/preprocessing.py| and \verb|data/README.md|.

% \subsection{Required Actions}
% \subsubsection{Landing Page}
% The landing page presents a hero banner of Hong Kong's skyline and three feature cards that deep-link into the app: (i) \emph{Locations list}, (ii) \emph{Map view}, and (iii) \emph{Favourites}. The page is intentionally lightweight—no data fetch occurs here—so first paint is fast and navigation is obvious for new users.
% \paragraph{UI structure}
% \begin{itemize}
%   \item \textbf{Hero banner}: full-width image with subtle gradient overlay to preserve title contrast.
%   \item \textbf{Feature cards grid}: responsive 1–3 column layout with icon, title, and a one-line description; each card routes to its respective module.
%   \item \textbf{Top nav}: global navigation to Home, Location List, Event List, Map, and Favourite List; auth avatar with role badge.
% \end{itemize}
% \paragraph{Interactions}
% \begin{itemize}
%   \item Cards are keyboard-focusable and announce destination via \texttt{aria-label}.
%   \item No API calls are made on Home; all data work happens after navigation.
% \end{itemize}

% \subsubsection{Locations}
% \label{subsec:locations}
% This is a dedicated directory distinct from the Home page. It lists venues with sorting, text filters, district selection, distance range (if geolocation granted), and a favourites toggle.
% CRUD for Admin Users: Administrators can create, edit, and delete locations. The create/edit sheets include name, district, latitude, and longitude, accompanied by a live map preview to verify coordinates. The list view supports filtering by name/district and, where geolocation is granted, a computed distance column. Successful mutations refresh the table so changes appear without reloading.
% \paragraph{Key features}
% \begin{itemize}
%   \item \textbf{Client-side filters}: name/description/district; range slider binds to a custom filter predicate on the distance column.
%   \item \textbf{Distance computation}: when permitted, user coordinates feed a Haversine helper to derive per-row distance (km).
%   \item \textbf{Favourites}: star button writes to \texttt{localStorage} (\texttt{favourites/v1}) with optimistic UI.
%   \item \textbf{Admin actions}: left-sheet modals for \emph{Create} and \emph{Edit}; deletion with confirm dialog; table refreshes on success.
% \end{itemize}
% \paragraph{CRUD Admin}
% \label{subsec:locations-crud}
% Parallel CRUD to events:
% \begin{itemize}
%   \item \textbf{Create/Edit sheets}: fields for name, district, latitude, longitude; live map marker preview.
%   \item \textbf{Success flow}: close sheet, toast message, and in-memory table refresh (\emph{no} hard reload required).
% \end{itemize}

% \subsubsection{Events}
% CRUD for Admin Users: Administrators can list, search, sort, create, edit, and delete events. The table includes fixed column widths with truncation to prevent long content from collapsing other columns, and a role-gated toolbar exposes action buttons only to admins. Creating or editing opens a Radix Sheet containing form fields for title, description, date, time, price, presenters, and venue (select). Required fields are validated on submit. After a successful mutation, the table refreshes and a “Last Updated” indicator (persisted in \texttt{localStorage}) is shown in the toolbar.
% \paragraph{Create Event}
% \begin{itemize}
%   \item Fields: title, description, venue (select), date, time, price, presenters; embedded map previews venue coordinates.
%   \item On submit, payload validation occurs client-side; successful POST closes the sheet and triggers a table refresh.
% \end{itemize}
% \paragraph{Edit Event}
% \begin{itemize}
%   \item Opens with current values populated; allows partial edits.
%   \item Uses \texttt{PUT /api/events/:id} (or \texttt{PATCH} if supported) and updates the visible row without a full page reload.
% \end{itemize}

% \subsubsection{Favourites View}
% The Favourites view reuses the table in Locations page, filtered to locations starred by the user. Favourites are persisted in \texttt{localStorage} using a simple key-value store (\texttt{favourites/v1}). Toggling the star updates the store immediately and re-renders the table without a server round-trip. The same filters/sorting and detail sheet are available.

% \subsubsection{API Integration and Auth}
% Frontend API helpers wrap REST endpoints for events and locations:
% \begin{itemize}
%   \item \textbf{Events:} \texttt{GET /api/events}, \texttt{POST /api/events}, \texttt{PUT /api/events/:id}, \texttt{DELETE /api/events/:id}.
%   \item \textbf{Locations:} \texttt{GET /api/locations}, \texttt{POST /api/locations}, \texttt{PUT /api/locations/:id}, \texttt{DELETE /api/locations/:id}.
% \end{itemize}
% Requests include \texttt{Authorization: Bearer \textless token\textgreater} retrieved from \texttt{localStorage}. Responses are normalized for table display and for form initial values in the edit sheets.

% \subsection{Data Normalization \& UI Pre-processing}
% \textbf{Events:} The payload is normalized to a table row structure and a display-only \texttt{dateTime} string is composed from \texttt{date} and \texttt{time}. The presenters field is unified across backends emitting either \texttt{presenter} or \texttt{presenters}. For submissions, both \texttt{venueId} and a denormalized \texttt{venue} name are sent as required by the backend.
% \newline
% \textbf{Locations:} The API response is mapped to \{\texttt{id}, \texttt{name}, \texttt{district}, \texttt{latitude}, \texttt{longitude}, \texttt{num\_events}\}. When geolocation is available, a \texttt{distance} (km) field is computed for sorting and range filtering.

% \subsection{UI/UX Details}
% \begin{itemize}
%   \item \textbf{Fixed-width columns \& truncation} ensure long text does not hide other columns and maintain consistent row height.
%   \item \textbf{Accessible Sheets (Radix)} for Create/Edit: titles are provided; description/aria attributes are handled to suppress a11y warnings.
%   \item \textbf{Role-gated controls}: admin toolbars and destructive actions are hidden from non-admin users.
%   \item \textbf{Immediate feedback}: sheets close and tables refresh on success; validation errors are surfaced inline; a last-updated timestamp appears after any mutation.
% \end{itemize}

% \subsection{Minimal Backend Touchpoints for These Features}
% To support the above flows, the Events and Locations routes accept create/update/delete with the field names and payload shapes described. Events update requires both \texttt{venueId} and denormalized \texttt{venue}; presenters are handled when the backend emits either \texttt{presenter} or \texttt{presenters}. Responses return the normalized fields consumed by the table and edit sheets.

% \subsection{Extra Features}

% \subsection{Algorithms/Functions}
% \begin{itemize}
%     \item \textbf{Haversine distance}: computes user-to-venue distance for sorting and range filtering.
%     \item \textbf{Debounced filters}: search inputs debounce updates to avoid excessive re-rendering.
%     \item \textbf{Fetch wrappers}: centralized REST helpers attach \texttt{Authorization} headers and handle JSON/HTTP errors.
%     \item \textbf{Optimistic UI}: favourites toggle updates the view immediately; state reconciles on storage events.
% \end{itemize}

% \subsection{Data Schemas and Models}
% \textbf{Event} (consumed fields): \texttt{\_id}, \texttt{title}, \texttt{description}, \texttt{date}, \texttt{time}, \texttt{price}, \texttt{presenter/presenters}, \texttt{venueId}, \texttt{venue}.
% \newline
% \textbf{Location} (consumed fields): \texttt{\_id} (\(\rightarrow\) \texttt{id}), \texttt{nameE} (\(\rightarrow\) \texttt{name}), \texttt{district}, \texttt{latitude}, \texttt{longitude}, \texttt{num\_events}.
% \newline
% Client-side validation ensures Event submissions include \texttt{title}, \texttt{date}, \texttt{time}, \texttt{venueId}, and \texttt{venue}; and Location submissions include \texttt{nameE}, \texttt{district}, \texttt{latitude}, and \texttt{longitude}.
